\section{How to Read a Field Equation}

A field equation should first be read as a local statement. It relates the value
of a field and its derivatives at a point, or in an arbitrarily small
neighborhood of that point. This locality is one of the defining ideas of field
theory. What happens here is controlled by the field here and by how it changes
nearby.

For example, an equation such as
\[
    (\Box+m^2)\phi=0
\]
is not merely a compact formula. It says that time variation, spatial variation,
and the local value of the field balance one another. When \(m=0\), the equation
has the character of wave propagation. When \(m\neq0\), the field also has a
local restoring term. The parameter \(m\) therefore changes the physical behavior
of solutions.

Whenever a new field equation appears, the following questions should be asked
before doing calculations:
\begin{enumerate}
    \item What is the unknown field, and what kind of object is it?
    \item What is the domain: space, spacetime, or a curved manifold?
    \item Which derivatives appear, and what local behavior do they measure?
    \item What constants occur, and what are their dimensions?
    \item What initial or boundary data are needed?
    \item Does the equation come from an action principle?
    \item Which symmetries and conservation laws does the equation express?
    \item Which quantities are gauge-invariant or physically observable?
\end{enumerate}

This habit prevents a common mistake: treating equations as isolated formulas.
In physics, an equation is compressed information about a model. It must be
unpacked.

\begin{examplebox}
Maxwell's equations are not only equations for \(\vect{E}\) and \(\vect{B}\).
They encode charge conservation, propagation of electromagnetic waves, gauge
redundancy, local energy flow, and the coupling between fields and sources.
\end{examplebox}

A good field theorist does not only solve equations. They ask what the equation
is allowed to mean. What has been idealized? What is held fixed? What counts as a
source? Which transformations change the physical state, and which merely change
the description? These questions are as important as the algebra.
